%% kapa: exam=1 round=20 year=2009 period=1 form=A subject=civil_law subject_ko=민법 seq=36 number=36
%% source: Q-Net 감정평가사 기출문제 (cst00309) artlSeq=5024178 file=extracted/1008/5024178/2009년 제20회 감정평가사 1차시험 기출문제/1교시.hwp
%% conversion: kordoc 4.24.51 -> tools/gen1.py (auto); review=proofread-sonnet; flags=none
\begin{question}{36}
저당권의 효력이 미치는 범위에 관한 설명으로 옳지 \uline{않은} 것은?
\begin{choices}
  \item 저당권 설정 후 건물을 증축하여 증축 부분이 기존 저당 목적 건물에 부합한 경우 그 부합된 부분에도 저당권의 효력이 미친다.
  \item 타인이 저당목적물의 상용에 제공한 물건에 대해서는 저당권의 효력이 미치지 않는다.
  \item 건물에 대한 저당권의 효력은 그 건물의 소유를 목적으로 하는 지상권에도 미친다.
  \item 저당목적물이 매매된 경우 그 매매대금에 대하여도 저당권자가 물상대위할 수 있다.
  \item 저당권자가 물상대위권을 행사하기 위한 요건으로서의 지급 또는 인도 전 압류는 반드시 저당권자가 할 필요는 없다.
\end{choices}
\end{question}
\answer{4}
